\SourceLeaf{171}%
% A544-P0171-D01
nicht ein besonderer Umstand ihn dazu veranlasst hätte, der hier erwähnt werden muss. Seit dem Jahre 1805 hatte er nemlich angefangen, sich mit den Theorien der cubischen und biquadratischen Reste zu beschäftigen, welche noch weit reichhaltiger und interessanter sind, als die Theorie der quadratischen Reste. Es zeigten sich bei jenen Untersuchungen dieselben Erscheinungen wie bei der letztern, nur gleichsam mit vergrössertem Massstabe. Durch Induction, sobald nur der rechte Weg dazu eingeschlagen war, fanden sich sogleich eine Anzahl höchst einfacher Theoreme, die jene Theorien ganz erschöpfen, mit den für die quadratischen Reste geltenden Lehrsätzen eine überraschende Aehnlichkeit haben, und namentlich auch zu dem Fundamentaltheorem das Gegenstück darbieten. Allein die Schwierigkeiten, für jene Lehrsätze ganz befriedigende Beweise zu finden, zeigten sich hier noch viel grösser, und erst nach vielen, eine ziemliche Reihe von Jahren hindurch fortgesetzten Versuchen ist es dem Verfasser endlich gelungen, sein Ziel zu erreichen. Die grosse Analogie der Lehrsätze selbst, bei den quadratischen und bei den höhern Resten, liess vermuthen, dass es auch analoge Beweise für jene und diese geben müsse; allein die zuerst für die quadratischen Reste gefundenen Beweisarten vertrugen gar keine Anwendung auf die höhern Reste, und gerade dieser Umstand war der Beweggrund, für jene immer noch andere neue Beweise aufzusuchen. Der Verf. wünscht daher, dass man die vorliegende Abhandlung, die für die Theorie der quadratischen Reste noch einige neue Hülfsquellen eröffnet, als Vorläuferin der Theorie der cubischen und biquadratischen Reste betrachte, die er in Zukunft bekannt zu machen denkt, und die zu den schwierigsten Gegenständen der höhern Arithmetik gehören.

% A544-P0171-D02
Die gegenwärtige Abhandlung besteht aus dreien von einander unabhängigen Theilen. Sie enthält nemlich den fünften und sechsten Beweis des Fundamental-Theorems und eine neue, mit dem dritten Beweise zusammenhängende Methode,\textperiodcentered\ zu entscheiden, ob eine vorgegebene ganze Zahl von einer gegebenen Primzahl quadratischer Rest oder Nichtrest sei. Unter den vier ersten Beweisen war der dritte unstreitig derjenige, der die grösste Einfachheit mit Unabhängigkeit von fremdartigen Untersuchungen vereinigte, daher ihn auch \textsc{Legendre} in die neue Ausgabe seines \emph{Essai d'une théorie des nombres} aufgenommen hat. Der \emph{fünfte} Beweis scheint dem dritten in beiden Hinsichten wenigstens gleich zu kommen. Beide Beweise haben insofern einige Verwandtschaft, dass sie von einem und demselben Lehrsatze ausgehen, sind aber bei der weitern Ausführung völlig von \SourceLeaf{172}%
% A544-P0172-D01
einander verschieden. Dieser Lehrsatz besteht in Folgendem: Wenn $m$ eine (positive ungerade) Primzahl; $M$ eine ganze durch $m$ nicht theilbare Zahl bedeutet, wenn ferner unter den Resten, die aus der Division der Producte
\[
M,\ 2M,\ 3M,\ 4M\mathinner{.\,.\,.\,.\,.\,.}\tfrac12(m-1)M
\]
durch $m$ entstehen, die Anzahl derjenigen, die grösser als $\frac12m$ sind, durch $n$ bezeichnet wird, so ist $M$ quadratischer Rest oder Nichtrest von $m$, jenachdem $n$ gerade oder ungerade ist. Um nun zu dem Beweise des Fundamentallehrsatzes zu gelangen, wird angenommen, dass auch $M$ eine ungerade positive Primzahl und $N$ in Beziehung auf $M$ und $m$ dasselbe bedeutet, was $n$ in Beziehung auf $m$ und $M$ ausdrückt, so dass $N$ gerade oder ungerade entscheidet, ob $m$ quadratischer Rest oder Nichtrest von $M$ ist. Durch eine sehr kurze Reihe von Schlüssen zeigt der Verfasser, dass die Anzahl aller positiven ganzen Zahlen, die zugleich kleiner als $\frac12mM$ sind, mit $m$ dividirt einen Rest kleiner als $\frac12m$, und mit $M$ dividirt einen Rest kleiner als $\frac12M$ geben,
\[
=\tfrac18(m-1)(M-1)+\tfrac12n+\tfrac12N
\]
und folglich allemal
\[
\tfrac14(m-1)(M-1)+n+N
\]
eine gerade Zahl sei. So oft also wenigstens eine der Zahlen $m,M$ von der Form $4k+1$ ist, mithin $\frac14(m-1)(M-1)$ gerade, wird auch $n+N$ gerade sein, folglich entweder $n$ und $N$ beide gerade, oder beide ungerade. Wenn hingegen sowohl $m$ als $M$ von der Form $4k+3$ ist, wird nothwendig $n+N$ ungerade, folglich eine der Zahlen $n,N$ gerade, die andere ungerade sein. Hieraus folgt in Verbindung mit obigem Lehrsatze das Fundamental-Theorem von selbst.

% A544-P0172-D02
Der \emph{sechste} Beweis ist zwar von gleicher Kürze und Concinnität wie der fünfte, beruhet aber doch auf etwas künstlichern Combinationen. Der beschränkte Raum dieser Blätter erlaubt nur, mit Uebergehung des Einzelnen, hier das Hauptmoment zu berühren. Es bezeichnen
\begin{description}
\item[$p,q$] zwei (ungleiche positive ungerade) Primzahlen,
\item[$\alpha$] eine sogenannte \emph{radix primitiva} für den Modulus $p$, d. i. eine durch $p$ nicht theilbare (hier positive) ganze Zahl von der Art, dass keine niedrigere Potenz als $\alpha^{p-1}$ nach dem Modulus $p$ der Einheit congruent wird
\item[$x$] eine unbestimmte Grösse
\SourceLeaf{173}%
% A544-P0173-D01
\item[$\xi$] die Function
\[
x-x^\alpha+x^\zeta-x^\eta+x^\theta-\text{etc.}-x^\lambda
\]
wo (des bequemern Drucks wegen) $\zeta,\eta,\theta\ldots\lambda$ statt der Zahlen $\alpha\alpha,\allowbreak \alpha^3,\allowbreak \alpha^4\ldots\allowbreak \alpha^{p-2}$ gesetzt sind;
\item[$\varepsilon$] die Einheit, positiv genommen, wenn $p$ von der Form $4k+1$, negativ, wenn $p$ von der Form $4k+3$ ist;
\item[$\delta$] die Einheit, positiv genommen, wenn wenigstens eine der Zahlen $p,q$ von der Form $4k+1$ ist, negativ, wenn beide von der Form $4k+3$ sind;
\item[$\gamma$] die Einheit, positiv genommen, wenn $q$ ein quadratischer Rest von $p$ ist, negativ, wenn $q$ quadratischer Nichtrest von $p$ ist;
\item[$\beta$] die Einheit, positiv genommen, wenn $p$ ein quadratischer Rest von $q$, negativ, wenn $p$ quadratischer Nichtrest von $q$ ist.
\end{description}

% A544-P0173-D02
Nach diesen Vorbereitungen folgt leicht aus dem 51. Art. der \emph{Disquisitiones Arithmeticae}, dass die Function
\[
\xi^q-x^q+x^{q\alpha}-x^{q\zeta}+x^{q\eta}-x^{q\theta}+\text{etc.}+x^{q\lambda}
\]
entwickelt lauter durch $q$ theilbare Coëfficienten bekommt, und daher, wenn diese Function $=qX$ gesetzt wird, $X$ eine auch in Beziehung auf die Coëfficienten \emph{ganze} Function werde. Durch Schlüsse, in die näher einzugehen hier zu weitläufig sein würde, wird in der Abhandlung bewiesen, dass die Function $qX\xi$ mit $x^{p-1}+x^{p-2}+x^{p-3}+x^{p-4}+\text{etc.}+x+1$ dividirt, den Rest
\[
\varepsilon p(\delta p^{\frac12(q-1)}-\gamma)
\]
gibt, daher aus der Division der Function $X\xi$ mit demselben Divisor der Rest
\[
\frac{\varepsilon p(\delta p^{\frac12(q-1)}-\gamma)}{q}
\]
hervorgehen wird. Diese Grösse muss daher nothwendig eine ganze Zahl sein, woraus, weil $\delta\delta=1$ ist, leicht geschlossen wird, dass
\[
p^{\frac12(q-1)}-\gamma\delta
\]
durch $q$ theilbar sein müsse. Da nun auch $p^{\frac12(q-1)}-\beta$ durch $q$ nach einem bekannten Theorem theilbar ist, so wird nothwendig $\beta=\gamma\delta$ sein, woraus wiederum das Fundamental-Theorem von selbst folgt.

\SourceLeaf{174}%
% A544-P0174-D01
Das Fundamental-Theorem, verbunden mit einigen bekannten Lehrsätzen, kann zwar zu einer ziemlich kurzen Auflösung der Aufgabe dienen, zu entscheiden, ob eine vorgegebene ganze positive Zahl von einer gegebenen Primzahl quadratischer Rest oder Nichtrest sei, wie in der Abhandlung ausführlich gezeigt ist. Allein bei weiterm Nachdenken über den dritten Beweis des Fundamental-Theorems kam der Verf. auf eine noch viel geschmeidigere Auflösung, welche die dritte Abtheilung der Abhandlung ausmacht, und wovon wir hier blos die Endregel hersetzen, indem wir die Entwickelung ihrer Gründe Kürze halber übergehen. Wenn entschieden werden soll, ob die ganze positive Zahl $b$, welche durch die Primzahl $a$ nicht theilbar ist, von dieser ein quadratischer Rest oder Nichtrest sei, so bilde man, ganz auf dieselbe Art, wie wenn der grösste gemeinschaftliche Divisor von $a$ und $b$ gesucht werden sollte, die Gleichungen
\[
\begin{aligned}
a&=\beta b+c\\
b&=\gamma c+d\\
c&=\delta d+e\\
d&=\varepsilon e+f\text{ u. s. w.}
\end{aligned}
\]
bis man in der Reihe der Zahlen $a,b,c,d,e,f$ u. s. w. auf .die Einheit kommt. Man. bezeichne die Zahlen $\frac12a,\frac12b,\frac12c,\frac12d$ u. s. w., mit Weglassung des ihnen anhängenden Bruches $\frac12$, in so fern einige der Zahlen $a,b,c,d$ u. s. w. ungerade sind, durch $a',b',c',d'$ u. s. w.; man nenne $\mu$ die Anzahl der in der Reihe $a',b',c',d'$ u. s. w. vorkommenden Folgen zweier ungeraden Zahlen unmittelbar nach einander, endlich nenne man $\nu$ die Anzahl derjenigen ungeraden Zahlen in der Reihe $\beta,\gamma,\delta,\varepsilon$ u. s. w., welchen in der Reihe $b',c',d',e'$ u. s. w. der Ordnung nach eine Zahl von der Form $4k+1$ oder $4k+2$ entspricht. Dieses vorausgesetzt, wird $b$ quadratischer Rest oder Nichtrest von $a$ sein, je nachdem $\mu+\nu$ gerade oder ungerade ist, den einzigen Fall ausgenommen, wo zugleich $b$ gerade und $a$ von der Form $8k+3$ oder $8k+5$ ist, in welchem von jener Regel das Gegentheil Statt findet, so dass ein gerades $\mu+\nu$ anzeigt, dass $b$ quadratischer Nichtrest von $a$ ist, ein ungerades $\mu+\nu$ hingegen, dass $b$ quadratischer Rest von $a$ ist.
\BookRule{.28\linewidth}

\clearpage
\SourceLeaf{175}%
% A544-P0175-D01
\pdfbookmark[0]{Göttingische gelehrta Anzeigen. 1825 April 11.}{notice1825april11}
\begin{center}
\rule{.82\linewidth}{.35pt}\par
.Göttingische gelehrta Anzeigen.\quad 1825 April 11.\par
\rule{.82\linewidth}{.35pt}
\end{center}
Am 5. April überreichte Hr. Hofr. \textsc{Gauss} der Königl. Societät eine Vorlesung, überschrieben:
\begin{center}\itshape
Theoria Residuorum Biquadraticorum, Commentatio prima.
\end{center}

% A544-P0175-D02
Die Theorie der quadratischen Reste bildet bekanntlich einen der interessantesten Theile der Höhern Arithmetik, welchen man jetzt nach vielfach wiederholten Untersuchungen als vollendet und abgeschlossen betrachten kann: die Geschichte desselben betreffende Nachrichten findet man in diesen Blättern 1808 Mai 12 und Sept. 19, und 1817 März 10. An letzterm Orte sind auch bereits einige vorläufige Nachrichten über die Nachforschungen mitgetheilt, welche der Verfasser der vorliegenden Abhandlung seit dem Jahre 1805 über die verwandte, eben so fruchtbare und interessante, aber sehr viel schwierigere Theorie der cubischen und biquadratischen Reste angestellt hatte. Obgleich schon damals im Besitz der wesentlichen Momente dieser Theorie, ist er doch bisher durch andere Arbeiten abgehalten, öffentlich etwas davon bekannt zu machen, und erst jetzt ist es ihm möglich geworden, sich mit der Ausarbeitung eines Theils dieser Untersuchungen zu beschäftigen. Der Anfang ist jetzt mit der Theorie der biquadratischen Reste gemacht, die der Theorie der quadratischen Reste näher verwandt ist, als die der cubischen. Inzwischen ist die gegenwärtige Abhandlung
